\documentclass[11pt,a4paper]{article}

% ── Packages ──────────────────────────────────────────────
\usepackage[utf8]{inputenc}
\usepackage[T1]{fontenc}
\usepackage[margin=2.2cm]{geometry}
\usepackage{xcolor}
\usepackage{tcolorbox}
\usepackage{hyperref}
\usepackage{enumitem}
\usepackage{tabularx}
\usepackage{booktabs}
\usepackage{fancyhdr}
\usepackage{parskip}

% ── Colors ────────────────────────────────────────────────
\definecolor{radioBlue}{HTML}{1a5276}
\definecolor{radioCyan}{HTML}{17a2b8}
\definecolor{radioGreen}{HTML}{28a745}
\definecolor{radioYellow}{HTML}{e6a817}
\definecolor{radioRed}{HTML}{dc3545}
\definecolor{radioDark}{HTML}{212529}
\definecolor{radioLight}{HTML}{f8f9fa}
\definecolor{radioGray}{HTML}{6c757d}

% ── Hyperlinks ────────────────────────────────────────────
\hypersetup{
    colorlinks=true,
    linkcolor=radioBlue,
    urlcolor=radioCyan,
    pdftitle={RadioSim User Guide},
    pdfauthor={RadioSim},
    pdfsubject={Radio Stations Simulation Tool}
}

% ── Headers & Footers ─────────────────────────────────────
\pagestyle{fancy}
\fancyhf{}
\fancyhead[L]{\textcolor{radioGray}{\small RadioSim}}
\fancyhead[R]{\textcolor{radioGray}{\small User Guide v1.0}}
\fancyfoot[C]{\textcolor{radioGray}{\small --- \thepage{} ---}}
\setlength{\headheight}{14.5pt}
\renewcommand{\headrulewidth}{0.4pt}
\renewcommand{\footrulewidth}{0.4pt}

% ── Custom boxes ──────────────────────────────────────────
\newtcolorbox{tipbox}{
    colback=radioLight!50, colframe=radioGreen,
    left=2mm, right=2mm, top=2mm, bottom=2mm,
    before upper={\textbf{\textcolor{radioGreen}{[i]}}\enspace}
}

\newtcolorbox{warnbox}{
    colback=radioYellow!15, colframe=radioYellow,
    left=2mm, right=2mm, top=2mm, bottom=2mm,
    before upper={\textbf{\textcolor{radioYellow}{[!]}}\enspace}
}

% ── Commands ──────────────────────────────────────────────
\newcommand{\cmd}[1]{\textcolor{radioDark}{\texttt{#1}}}
\newcommand{\key}[1]{\textcolor{radioBlue}{\framebox{\texttt{#1}}}}
\newcommand{\signalbar}[1]{%
    {\color{radioGreen}\rule{#1mm}{3mm}}%
}

% ── Title ─────────────────────────────────────────────────
\title{\vspace{-1cm}\Huge\textcolor{radioBlue}{\textbf{RadioSim}}\\[2mm]
    \Large Radio Stations Simulation Tool\\[2mm]
    \normalsize\textcolor{radioGray}{User Guide {\tiny —} v1.0}}
\author{}
\date{}

\begin{document}

\maketitle
\thispagestyle{fancy}

\begin{center}
\begin{tcolorbox}[colback=radioBlue!8, colframe=radioBlue, width=0.92\textwidth,
    left=4mm, right=4mm, top=3mm, bottom=3mm]
    \textbf{What is RadioSim?}\quad
    RadioSim lets you \textbf{listen to your music as if it were broadcast over the radio}.
    Feed it your MP3 files or YouTube links, pick a radio mode (FM, AM, DAB+\dots),
    and hear what it would sound like over the air --- complete with static,
    signal fading, and all the quirks of real radio reception.
\end{tcolorbox}
\end{center}

% ═══════════════════════════════════════════════════════════
\section{Quick Start}
% ═══════════════════════════════════════════════════════════

You need \textbf{Python 3.10+} and \textbf{ffmpeg}. That's it.

\begin{enumerate}[leftmargin=*]
    \item Open a terminal in the \texttt{radio-sim/} folder.
    \item Activate the environment:
    \begin{center}
        \cmd{source venv/bin/activate}
    \end{center}
    \item Play your music folder as an FM station:
    \begin{center}
        \cmd{python -m src.main --mode fm --freq 101.1 --source \textasciitilde/Music/}
    \end{center}
\end{enumerate}

\begin{tipbox}
    \textbf{First time?} Run \cmd{pip install -r requirements.txt} once to install dependencies.
\end{tipbox}

\vspace{3mm}
\textbf{Three things to try right now:}

\begin{tcolorbox}[colback=radioGreen!8, colframe=radioGreen, left=4mm, right=4mm]
\begin{enumerate}[nosep]
    \item \cmd{python -m src.main --mode fm  --freq 101.1 --source \textasciitilde/Music/}
    \item \cmd{python -m src.main --mode am  --freq 880  --source \textasciitilde/Music/ --rssi -80}
    \item \cmd{python -m src.main --mode dab --freq 202.0 --source \textasciitilde/Music/}
\end{enumerate}
\end{tcolorbox}

% ═══════════════════════════════════════════════════════════
\section{The Five Radio Modes}
% ═══════════════════════════════════════════════════════════

RadioSim simulates five broadcast technologies. Each one sounds different.
Here's what to expect.

% ── FM ──
\subsection*{\textcolor{radioBlue}{\textbf{FM} --- Frequency Modulation}}

\begin{tabularx}{\textwidth}{@{}>{\raggedright\arraybackslash}p{4.5cm} X@{}}
    \textbf{Frequency band}  & 87.5 -- 108.0\,MHz \\
    \textbf{Audio quality}   & 50\,Hz -- 15\,kHz, stereo \\
    \textbf{Sound character} & Clean with slight background hiss. \\
                             & At weak signal: echoes (multipath), stereo blends to mono. \\
    \textbf{Tuning step}     & 0.1\,MHz \\
\end{tabularx}

\vspace{2mm}
\begin{tabular}{ccl}
    \toprule
    \textbf{Signal} & \textbf{RSSI} & \textbf{What you hear} \\
    \midrule
    \signalbar{13} Excellent  & $-30$ dBm & Studio quality, near-perfect \\
    \signalbar{10} Very good  & $-45$ dBm & Clean, barely audible hiss \\
    \signalbar{7}  Fair       & $-60$ dBm & Noticeable hiss, slight fading \\
    \signalbar{4}  Weak       & $-75$ dBm & Heavy hiss, forced mono \\
    \signalbar{1}  Threshold  & $-90$ dBm & Noise dominates the signal \\
    \bottomrule
\end{tabular}

% ── AM ──
\subsection*{\textcolor{radioRed}{\textbf{AM} --- Amplitude Modulation}}

\begin{tabularx}{\textwidth}{@{}>{\raggedright\arraybackslash}p{4.5cm} X@{}}
    \textbf{Frequency band}  & 530 -- 1710\,kHz \\
    \textbf{Audio quality}   & 50\,Hz -- 5\,kHz, mono only \\
    \textbf{Sound character} & ``Frying'' static, slow fading, crackle from lightning. \\
                             & Narrow bandwidth --- sounds like a telephone line. \\
    \textbf{Tuning step}     & 1\,kHz \\
\end{tabularx}
\vspace{1mm}

\begin{warnbox}
    AM radio is \textbf{always mono} and \textbf{always noisy}.
    Try it at \cmd{--rssi -80} for the full ``distant station at midnight'' experience.
\end{warnbox}

% ── AMHD ──
\subsection*{\textcolor{radioYellow}{\textbf{AMHD} --- AM HD Radio}}
Hybrid digital AM. Same dial as AM (530--1710\,kHz), but a digital sidechannel
adds stereo and extended bandwidth (15\,kHz) at strong signal. Falls back to
analog AM quality when the signal weakens. The best of both worlds --- when
you're close enough to the tower.

% ── FMHD ──
\subsection*{\textcolor{radioGreen}{\textbf{FMHD} --- FM HD Radio}}
Hybrid digital FM. Full CD bandwidth (20\,Hz -- 20\,kHz), discrete stereo.
Supports multicasting: up to 3 sub-channels on one frequency. Near-transparent
at strong signal. Almost indistinguishable from the original file.

% ── DAB+ ──
\subsection*{\textcolor{radioCyan}{\textbf{DAB+} --- Digital Audio Broadcasting}}

\begin{tabularx}{\textwidth}{@{}>{\raggedright\arraybackslash}p{4.5cm} X@{}}
    \textbf{Frequency band}  & 174 -- 240\,MHz \\
    \textbf{Audio quality}   & 20\,Hz -- 20\,kHz, stereo, HE-AAC \\
    \textbf{Sound character} & Pristine digital audio --- \textbf{or total silence}. \\
                             & No static. No hiss. No gradual fade. \\
                             & The ``cliff effect'': perfect above threshold, mute below. \\
\end{tabularx}

\vspace{2mm}
\begin{center}
\textit{``DAB doesn't get noisy. It just stops.''}
\end{center}

\vspace{2mm}

\begin{center}
\begin{tcolorbox}[colback=radioLight, colframe=radioGray, width=0.85\textwidth]
\textbf{Quick Comparison Table}\\[2mm]
\begin{tabular}{@{}lcccc@{}}
    \toprule
    \textbf{Feature} & \textbf{FM} & \textbf{AM} & \textbf{AMHD} & \textbf{FMHD} \\
    \midrule
    Bandwidth (max) & 15\,kHz & 5\,kHz & 15\,kHz & 20\,kHz \\
    Stereo          & Yes & No & Yes (digital) & Yes \\
    Noise type      & Hiss & Static & Digital artifacts & Near-none \\
    Gradual fade?   & Yes & Yes & Crossfade & Crossfade \\
    \bottomrule
\end{tabular}

\vspace{1mm}
\textbf{DAB+}: 20\,kHz, stereo, digital --- but mutes completely below threshold.
\end{tcolorbox}
\end{center}

% ═══════════════════════════════════════════════════════════
\section{Terminal UI --- Interactive Mode}
% ═══════════════════════════════════════════════════════════

Launch without \cmd{--no-tui} to get the interactive radio faceplate:

\begin{center}
\cmd{python -m src.main --mode fm --freq 101.1 --source \textasciitilde/Music/}
\end{center}

\subsection*{Keyboard Controls}

\begin{center}
\begin{tabular}{@{}ll@{\hspace{15mm}}ll@{}}
    \toprule
    \textbf{Key} & \textbf{Action} & \textbf{Key} & \textbf{Action} \\
    \midrule
    \key{Q} / Esc    & Quit             & \key{M} / Tab & Cycle mode \\
    \key{$\leftarrow$} \key{$\rightarrow$} & Tune frequency   & \key{$\uparrow$} \key{$\downarrow$} & Volume \\
    \key{+} \key{-}  & Volume (alt)     & \key{[} \key{]} & RSSI weaker/stronger \\
    \key{R}          & Randomize RSSI   & \key{N} & Next track \\
    \key{Space}      & Play / Pause     & \key{S} & Toggle source \\
    \bottomrule
\end{tabular}
\end{center}

\begin{tipbox}
    \textbf{Try this:} Launch FM at strong signal, then press \key{[} repeatedly.
    You'll hear the signal degrade in real time --- like driving away from a city.
\end{tipbox}

% ═══════════════════════════════════════════════════════════
\section{Command-Line Reference}
% ═══════════════════════════════════════════════════════════

\begin{tabularx}{\textwidth}{@{}>{\ttfamily\footnotesize}l X@{}}
    \toprule
    \normalfont\textbf{Option} & \textbf{What it does} \\
    \midrule
    -m, --mode    & Radio mode: \texttt{fm}, \texttt{am}, \texttt{amhd}, \texttt{fmhd}, \texttt{dab} \\
    -f, --freq    & Frequency to tune (auto-selects center of band if omitted) \\
    -s, --source  & Path to MP3 folder/file, or YouTube URL \\
    -r, --rssi    & Signal strength in dBm ($-120$ to $0$, default $-45$) \\
    -v, --volume  & Output volume $0.0$ -- $1.0$ (default $0.8$) \\
    -w, --wav     & Export processed audio to a \texttt{.wav} file \\
    -d, --duration& Duration in seconds for WAV export (default $10$) \\
    --no-tui      & Headless mode --- no interactive display \\
    --help        & Show full help \\
    \bottomrule
\end{tabularx}

\subsection*{WAV Export}

Capture processed audio to compare modes or share with others:

{\footnotesize
\begin{tcolorbox}[colback=radioLight, colframe=radioGray, left=2mm, right=2mm]
\texttt{\# Export 30 seconds of AM radio with heavy static}\\
\texttt{python -m src.main --mode am --freq 880 --rssi -85 \textbackslash}\\
\texttt{\hspace{4mm} --source \textasciitilde/Music/song.mp3 \textbackslash}\\
\texttt{\hspace{4mm} --wav /tmp/am\_weak.wav --duration 30 --no-tui}
\end{tcolorbox}
}

% ═══════════════════════════════════════════════════════════
\section{Audio Sources}
% ═══════════════════════════════════════════════════════════

\subsection*{Local Files}
\begin{itemize}[nosep]
    \item Point \cmd{--source} at a folder or a single file.
    \item Folders are scanned \textbf{recursively} by default.
    \item Files are \textbf{shuffled} --- each session is a different playlist.
    \item Supported formats: MP3, WAV, FLAC, OGG, M4A, AAC, Opus.
\end{itemize}

\subsection*{YouTube}
\begin{itemize}[nosep]
    \item Paste any YouTube URL as \cmd{--source}.
    \item Audio streams in real time --- no waiting for download.
    \item Requires \texttt{ffmpeg} on your system.
\end{itemize}

\begin{center}
\cmd{python -m src.main --mode am --freq 880 --source "https://www.youtube.com/watch?v=..."}
\end{center}

% ═══════════════════════════════════════════════════════════
\section{Recipes}
% ═══════════════════════════════════════════════════════════

\begin{tcolorbox}[colback=radioBlue!8, colframe=radioBlue,
    title={\textbf{Compare All 5 Modes (same song)}}]
{\footnotesize
\texttt{for mode in fm am amhd fmhd dab; do}\\
\texttt{\hspace{4mm}python -m src.main --mode \$mode --wav /tmp/cmp\_\$\{mode\}.wav}\\
\texttt{\hspace{8mm} --source song.mp3 --no-tui}\\
\texttt{done}
}
\end{tcolorbox}

\begin{tcolorbox}[colback=radioGreen!8, colframe=radioGreen,
    title={\textbf{Driving Away from a Station}}]
{\footnotesize
\texttt{python -m src.main --mode fm --freq 101.1 \textbackslash}\\
\texttt{\hspace{4mm} --source \textasciitilde/Music/ --rssi -30}\\[1mm]
Then press \key{[} repeatedly --- hear the signal get weaker.
}
\end{tcolorbox}

\begin{tcolorbox}[colback=radioYellow!8, colframe=radioYellow,
    title={\textbf{Distant AM Station at Night}}]
{\footnotesize
\texttt{python -m src.main --mode am --freq 880 --rssi -85 \textbackslash}\\
\texttt{\hspace{4mm} --source \textasciitilde/Music/}\\[1mm]
Heavy static, deep fading, occasional crackle. Pure ionosphere.
}
\end{tcolorbox}

\begin{tcolorbox}[colback=radioCyan!8, colframe=radioCyan,
    title={\textbf{DAB+ Cliff Effect Demo}}]
{\footnotesize
\texttt{python -m src.main --mode dab --freq 202.0 --rssi -85 \textbackslash}\\
\texttt{\hspace{4mm} --source \textasciitilde/Music/}\\[1mm]
Right on the digital cliff --- alternates between perfect and silence.
}
\end{tcolorbox}

% ═══════════════════════════════════════════════════════════
\section{Troubleshooting}
% ═══════════════════════════════════════════════════════════

\begin{tabularx}{\textwidth}{@{}>{\raggedright\arraybackslash}p{4cm} X@{}}
    \toprule
    \textbf{Problem} & \textbf{Solution} \\
    \midrule
    No sound &
    Run \texttt{aplay -l} to check audio hardware is detected. \\
    & Run \texttt{alsamixer} --- your output might be muted. \\
    \addlinespace
    \texttt{ffmpeg not found} &
    Install: \texttt{sudo apt install ffmpeg} (Debian) or \texttt{brew install ffmpeg} (macOS). \\
    \addlinespace
    Audio stutters / glitches &
    Edit \texttt{src/engine/pipeline.py}: increase \texttt{buffer\_chunks} from 32 to 64. \\
    \addlinespace
    Import errors &
    Always run from the \texttt{radio-sim/} directory with \texttt{source venv/bin/activate}. \\
    \addlinespace
    YouTube not working &
    Update yt-dlp: \texttt{pip install --upgrade yt-dlp} \\
    \addlinespace
    \texttt{venv/bin/activate} not found &
    Create it first: \texttt{python3 -m venv venv \&\& pip install -r requirements.txt} \\
    \bottomrule
\end{tabularx}

\vspace{5mm}
\begin{center}
\rule{0.5\textwidth}{0.4pt}\\[2mm]
{\small\textcolor{radioGray}{RadioSim {\tiny —} github.com/enigma {\tiny —} July 2026}}\\[1mm]
{\small\textcolor{radioGray}{Python {\tiny ·} SciPy {\tiny ·} Rich TUI {\tiny ·} ffmpeg {\tiny ·} yt-dlp}}
\end{center}

\end{document}
